Este capítulo sintetiza los hallazgos principales de la investigación, evalúa el cumplimiento de los objetivos planteados, expone las limitaciones del estudio, y; propone expectativas a futuro.

\section{Cumplimiento de Objetivos}

\subsection{Objetivo General}

El objetivo general era desarrollar una plataforma de inteligencia de negocio crediticia que, mediante un flujo de operaciones de Machine Learning automatizado, compare modelos de deep learning y gradient boosting para la predicción multi-horizonte de riesgo, e integre el modelo seleccionado en un datamart analítico con dashboards interactivos para el control de la mora.

\textbf{Este objetivo se cumplió de forma  satisfactoria}. Se implementó un sistema completo que incluye:

\begin{enumerate}
    \item Pipeline orquestado en Apache Airflow con dos flujos (entrenamiento e inferencia).
    \item Comparación de tres modelos: CNN, MLP y LightGBM. Monitoreado mediante MlFlow
    \item Selección de LightGBM como modelo ganador basado en métricas objetivas.
    \item Integración del modelo en un datamart analítico en esquema estrella.
    \item Dashboards interactivos en Apache Superset para visualización de resultados.
\end{enumerate}

\subsection{Objetivos Específicos}

\textbf{Objetivo 1: Identificar las causas del fracaso de las redes neuronales en la predicción de incumplimiento crediticio con datos tabulares.}

Los resultados son consistentes con Arram et al. \cite{arram2023credit}, quienes encontraron que MLP logra buen rendimiento solo con datasets grandes. En nuestro caso, con solo 293 muestras de entrenamiento, las redes neuronales no pudieron converger.

\textbf{Objetivo 2: Determinar qué características de los datos crediticios hacen que los modelos basados en árboles superen a las redes neuronales.}

Las características encontradas son:

\begin{enumerate}
    \item \textbf{Datos tabulares}: Las redes neuronales destacan en imágenes y texto, pero los árboles de decisión son superiores para datos estructurados en tabla.
    \item \textbf{Secuencias cortas}: La ventana de entrada (6 meses × 21 features = 128 valores) es insuficiente para que las CNN capturen patrones temporales relevantes.
    \item \textbf{Valores atípicos}: Los montos crediticios presentan valores extremos que afectan a las redes neuronales pero no a los árboles de decisión.
\end{enumerate}

\textbf{Objetivo 3: Evaluar el impacto del horizonte temporal en la degradación del rendimiento predictivo.}

Se confirmó que el rendimiento decae naturalmente con el horizonte:

\begin{itemize}
    \item \textbf{Horizontes 1-6 meses}: AUC-ROC > 83\%. Alta confiabilidad.
    \item \textbf{Horizontes 7-12 meses}: AUC-ROC 76-82\%. Confianza moderada.
    \item \textbf{Horizontes 13-18 meses}: AUC-ROC 70-77\%. Baja confianza.
\end{itemize}

Esta degradación se debe a la acumulación de incertidumbre en predicciones a mayor plazo y es consistente con la literatura existente.

\textbf{Objetivo 4: Visualizar un datamart analítico como parte de la infraestructura para integrar predicciones con datos históricos.}

Este objetivo se alinea con las recomendaciones de Kimball y Ross \cite{kimball2013datawarehouse} sobre el uso de esquemas estrella para datamarts analíticos. Se implementó un datamart en esquema estrella.

\section{Contribuciones}

Las principales contribuciones de este trabajo son:

\begin{enumerate}
    \item \textbf{Comparación rigurosa de modelos}: Se evaluaron tres enfoques de machine learning para predicción de riesgo crediticio multi-horizonte, demostrando que LightGBM supera a CNN y MLP en datos tabulares y con clases desbalanceadas. Estos resultados son consistentes con los hallazgos de Yu et al. \cite{yu2024advanced} y Zhu et al. \cite{zhu2024ensemble}.
    
    \item \textbf{Pipeline MLOps completo}: Se diseñó e implementó una sistema que cubre desde la ingreso de datos hasta la visualización, que automatizada en Airflow y tracking de experimentos en MLflow, siguiendo las prácticas de MLOps documentadas por Marcos-Mercadé et al. \cite{marcos2026empirical} y Micallef et al. \cite{micallef2026systematic}.
    
    \item \textbf{Datamart analítico integrado}: Se creó un esquema estrella que permite cruzar datos históricos con predicciones del modelo, facilitando el análisis exploratorio y la toma de decisiones.
    
    \item \textbf{Dashboards interactivos}: Se implementaron visualizaciones en Apache Superset que permiten a los usuarios de negocio explorar datos y predicciones.
\end{enumerate}

\section{Limitaciones del Estudio}

El presente trabajo presenta las siguientes limitaciones:

\begin{enumerate}
    \item \textbf{Degradación en horizontes largos}: Los horizontes 13-18 meses muestran una degradación significativa (AUC-ROC < 77\%), lo que reduce la utilidad de la predicción.
    
    \item \textbf{Ausencia de validación externa}: No se comparó con modelos de referencia externos ni se realizó validación cruzada con datos de otras instituciones, esto es entendible por la naturaleza restringida de la data.
    
    \item \textbf{Análisis de modelos mejor valorados}: No se exploraron técnicas avanzadas como embeddings de secuencias, attention mechanisms o transformers para datos tabulares, los cuales podrían mejorar el rendimiento de las redes neuronales según Nayak \cite{nayak2026calibrated}.
\end{enumerate}

\section{Perspectivas a Futuro}

Las líneas con visión de futura incluyen:

\begin{enumerate}
    \item \textbf{Aumentar el volumen de datos}: Incorporar datos de múltiples instituciones o generar datos sintéticos para mejorar el entrenamiento de redes neuronales. 
    
    \item \textbf{Optimizar horizontes largos}: Explorar técnicas de regularización, aprendizaje por transferencia o arquitecturas híbridas para mejorar el rendimiento en horizontes 13-18 meses.
    
    \item \textbf{Incorporar más fuentes de datos}: Integrar variables macroeconómicas, datos de crédito buró o información de redes sociales para enriquecer las predicciones.
    
    \item \textbf{Monitoreo en cambios en la data}: Implementar un sistema de alertas en el cambio de distribución de datos y active reentrenamiento automático.
    
    \item \textbf{Uso en giros de negocios semejantes}: El modelo de predicción se puede utilizar en otros giros de negocios que otorguen crédito directo como comercios de electrodomésticos, venta de autos, etc.
\end{enumerate}
